\documentclass[11pt]{article}

\usepackage[a4paper, margin=2.6cm]{geometry}
\usepackage{amsmath, amssymb}
\usepackage{graphicx}
\usepackage{booktabs}
\usepackage{caption}
\usepackage[hidelinks]{hyperref}
\usepackage{mathpazo}
\usepackage{microtype}

\captionsetup{font=small, labelfont=bf, margin=10pt}

\title{\textbf{Measuring Electromagnetic Induction: The Peak EMF Induced by a Magnet Falling Through a Coil}}
\author{Chhaya Ramnani (u8336188)\\ \small Lab partners: Tsz Wing Candice Ng, Chinatsu Komada\\ \small PHYS1201 --- Physics 2, Week 5 Laboratory: Electromagnetic Induction}
\date{Date of submission: 2026}

\begin{document}

\maketitle
\thispagestyle{empty}

\section*{1.0 Abstract}
This experiment examined how the peak electromotive force (emf) induced by a magnet falling through a coil depends on three things: 
the number of coil turns $N$, the coil radius $R$, and the magnet's speed $s$ as it passed through. The magnet's dipole moment was 
first found from calibration measurements of its on-axis field, giving $m = (1.18 \pm 0.07)$~A\,m$^2$. This was used to build a 
theoretical prediction for the peak emf with no free parameters, which the measured data from each of the three sub-experiments was 
then compared against directly --- no scaling factor was fitted anywhere. The peak emf rose roughly linearly with $N$ and with $s$, 
and fell off roughly as $R^{-2}$ with coil radius, matching the predicted shape $\varepsilon_{\text{peak}} \propto Ns/R^2$. None of 
the three datasets matched the absolute theoretical prediction within uncertainty, though, and the $N$ data disagreed by far the most 
(a measured-to-theory ratio of $0.461 \pm 0.003$). This was mainly put down to the magnet's finite size breaking the point-dipole 
assumption the theory relies on, and, specifically for the $N$ sub-experiment, to an unresolved mismatch between two setups that should 
have given the same reading.

\section*{2.0 Introduction}

\subsection*{2.1 Aims}
This experiment set out to see how the peak emf induced by a magnet falling through a coil depends on three parameters: the number of 
turns in the coil, $N$; the coil radius, $R$; and the magnet's speed as it passed through, $s$. A dipole model of the magnet was calibrated 
first from independent field measurements, and the resulting zero-parameter theoretical prediction was then compared directly against the 
measured emf for each sub-experiment, with no scaling factor fitted in.

The peak emf was expected to increase linearly with $N$, increase linearly with $s$, and fall off as $R^{-2}$ with coil radius, with all three trends together matching the combined proportionality $\varepsilon_{\text{peak}} \propto Ns/R^2$.

\subsection*{2.2 Background theory}
By Faraday's law, the emf induced in a circuit equals the negative rate of change of magnetic flux $\Phi_B$ through it:
\begin{equation}
\varepsilon = -\frac{d\Phi_B}{dt}, \qquad \Phi_B = \int B_\perp \, dA .
\end{equation}


\end{document}